% !PRINTED 705
\weberpage{705}
This function must therefore vanish in $\mathfrak P_2,\allowbreak \mathfrak P_3,\allowbreak \ldots$. At the point $\mathfrak P_2$, however, $z\,x_{1,1}\varrho_1,\allowbreak z\,x_{1,3}\varrho_3,\allowbreak \ldots,\allowbreak z\,x_{1,n}\varrho_n$ vanish, while $\varrho_2$ does not vanish. Consequently $z\,x_{1,2}$ must vanish in $\mathfrak P_2$, that is, for $z=\infty$, and $z^2 x_{1,2}$ must remain finite. We have therefore:

\gesperrt{If $r$ is not equal to $s$, then $z^2 x_{r,s}$ is finite for $z=\infty$.}

We now define the rational function $u^{(i)}$ by:
\weq{11}{x_{i,i} = a^{(i)}_{m-2} z^{m-2} + a^{(i)}_{m-3} z^{m-3} + \cdots + a^{(i)}_{-1} z^{-1} + u^{(i)} z^{-2},}
in which the constant $a^{(i)}_\nu$ is to be the same as in (5), and we obtain:
\weq{12}{\frac{dJ}{dz} - x_{i,i} = z^{-2}\bigl[\eta^{(i)} - u^{(i)}\bigr],}
and by (8):
\[
\bigl[\eta^{(i)} - u^{(i)}\bigr]\varrho_i = z^2 x_{i,1}\varrho_1 + \cdots + z^2 x_{i,i-1}\varrho_{i-1} + z^2 x_{i,i+1}\varrho_{i+1} + \cdots + z^2 x_{i,n}\varrho_n.
\]

At the point $\mathfrak P_i$, $\eta^{(i)}$ is finite and $\varrho_i$ different from zero, while on the right-hand side everything vanishes in $\mathfrak P_i$. Consequently $u^{(i)}$ too must be finite in $\mathfrak P_i$, and, since it is rational, also for $z=\infty$; and from (9) and (11) it follows that:
\weq{13}{S\left(\frac{dJ}{dz}\right) = \overset{i}{\Sigma} a^{(i)}_{m-2} z^{m-2} + \cdots + \overset{i}{\Sigma} a^{(i)}_{-1} z^{-1} + \overset{i}{\Sigma} u^{(i)} z^{-2}.}

By our assumption about $z$, the lower polygon $\mathfrak B$ of $dJ$ contains no point at which $z$ has a finite value.

If $\mathfrak P$ is a point at which $z$ has the finite value $z_0$, then
\weq{14}{(z - z_0) = \frac{\mathfrak N}{\mathfrak U}, \quad \frac{dJ}{dz} = \frac{\mathfrak A \mathfrak U^2}{\mathfrak Z\mathfrak B}, \quad (z-z_0)\frac{dJ}{dz} = \frac{\mathfrak U \mathfrak N \mathfrak A}{\mathfrak Z\mathfrak B},}
and the point $\mathfrak P$ is contained once more often in $\mathfrak N$ than in $\mathfrak Z$; consequently:
\weq{15}{\left[(z-z_0)\frac{dJ}{dz}\right]_0 = 0 \quad \text{(in $\mathfrak P$).}}
This is requirement a) in \S\,198.

If therefore $\lambda_1, \lambda_2, \ldots, \lambda_n$ is the normal basis in $z$, and $\mu_1, \mu_2, \ldots, \mu_n$ the complementary basis, then:
\weq{16}{\frac{dJ}{dz} = y_1 \mu_1 + y_2 \mu_2 + \cdots + y_n \mu_n,}
in which the $y_i$ are integral rational functions of $z$, from which it follows that
\weq{17}{S\left(\frac{dJ}{dz}\right) = y_1 \quad [\S\,198, (18)]}
\gesperrt{is an integral rational function of $z$.}
